% ---------- פרק 7: ניתוח סיכונים ותוכנית הפחתה ----------
% 22.11.2026: גוף שלישי, קיצור, איור בעברית. ערכי L ו-I הם הערכה איכותנית. ORG-04 (איש קשר בתע"א) הוסר לבקשת המרצה, 16.11.
\chapter{ניתוח סיכונים ותוכנית הפחתה}\label{ch:risks}

לכל סיכון מוערכות \textbf{הסתברות} $L$ ו\textbf{השפעה} $I$ בסולם 1 עד 5 (בהערכה איכותנית), וציון הסיכון הוא $E = L \times I$. הסימונים שונים מ-$P$ ו-$R$ של פרק~\ref{ch:design}. סיכון עם $E\ge 12$ הוא \textbf{קריטי}, ולכל סיכון נקבעות \textbf{פעולת מנע}, שננקטת מראש, ו\textbf{תוכנית מגירה}, שמופעלת אם הוא מתממש.

\section{זיהוי הסיכונים וכימותם}\label{sec:risk-id}

טבלה~\ref{tab:risks} מפרטת את הסיכונים, ואיור~\ref{fig:riskmatrix} ממקם אותם במטריצה. מבין ארבעת הסיכונים הקריטיים, שניים נוגעים לתוקף התיקוף (\en{DAT-01}, \en{VAL-02}) ושניים לאמינות הפלט של המערכת (\en{ALG-03}, \en{ENT-07}).

{\small
\begin{longtable}{|l|>{\raggedright\arraybackslash}p{7.2cm}|>{\raggedright\arraybackslash}p{1.9cm}|c|c|c|}
\caption{זיהוי הסיכונים וכימותם}\label{tab:risks}\\
\hline
\textbf{מזהה} & \textbf{תרחיש הכשל} & \textbf{תחום} & $L$ & $I$ & $E$ \\ \hline
\endfirsthead
\hline
\textbf{מזהה} & \textbf{תרחיש הכשל} & \textbf{תחום} & $L$ & $I$ & $E$ \\ \hline
\endhead
\en{DAT-01} & אין מספיק מידע ציבורי על מקרי התיקוף ועל אותות איכות, במיוחד בחברות פרטיות, ולכן אי אפשר לבדוק את \en{H1} ו-\en{H2} & נתונים & 4 & 4 & \textbf{16} \\ \hline
\en{VAL-02} & "ידיעת העתיד": המודל כבר יודע מה קרה לספק, והבדיקה לאחור נראית טובה מהמציאות & תיקוף & 4 & 4 & \textbf{16} \\ \hline
\en{ALG-03} & הזיה: הסוכן ממציא עובדה או מקור, והם נכנסים להמלצה & אלגוריתם & 3 & 5 & \textbf{15} \\ \hline
\en{SEC-04} & מידע פנימי נשלח בטעות לשירות ענן, או ששם הספק והרכיב חושפים כוונת רכש & אבטחת מידע & 2 & 5 & 10 \\ \hline
\en{LEG-05} & תנאי השימוש של מקור מידע אוסרים איסוף אוטומטי & משפטי & 3 & 3 & 9 \\ \hline
\en{CST-06} & עלות הקריאות למודל או מגבלות קצב חורגות מהתקציב & משאבים & 2 & 2 & 4 \\ \hline
\en{ENT-07} & זיהוי ישות שגוי: שם דומה, חברה-בת או שינוי שם גורמים לבדיקת הישות הלא נכונה, והשער מחמיץ ספק ברשימה & נתונים & 3 & 4 & \textbf{12} \\ \hline
\end{longtable}
}

\begin{figure}[!htbp]
\centering
\begin{tikzpicture}[x=0.95cm,y=0.6cm]
\foreach \p in {1,...,5}{
 \foreach \i in {1,...,5}{
   \pgfmathtruncatemacro{\r}{\p*\i}
   \ifnum\r>11 \def\col{red!35}\else\ifnum\r>5 \def\col{yellow!40}\else\def\col{green!20}\fi\fi
   \fill[\col] (\i-1,\p-1) rectangle (\i,\p);
   \draw[gray!50] (\i-1,\p-1) rectangle (\i,\p);
   \node[font=\tiny, gray] at (\i-0.15,\p-0.82) {\r};
 }}
\foreach \i in {1,...,5}{ \node[font=\scriptsize] at (\i-0.5,-0.3) {\i}; \node[font=\scriptsize] at (-0.3,\i-0.5) {\i}; }
\node[font=\small] at (2.5,-0.9) {\foreignlanguage{hebrew}{השפעה} $I$};
\node[font=\small, rotate=90] at (-0.9,2.5) {\foreignlanguage{hebrew}{הסתברות} $L$};
\node[font=\scriptsize, align=center] at (3.5,3.5) {DAT-01\\VAL-02};
\node[font=\scriptsize] at (4.5,2.5) {ALG-03};
\node[font=\scriptsize] at (3.5,2.5) {ENT-07};
\node[font=\scriptsize] at (4.5,1.5) {SEC-04};
\node[font=\scriptsize] at (2.5,2.5) {LEG-05};
\node[font=\scriptsize] at (1.5,1.5) {CST-06};
\end{tikzpicture}
\caption{מטריצת סיכונים: הסתברות מול השפעה. באדום הסיכונים הקריטיים ($E\ge 12$)}\label{fig:riskmatrix}
\end{figure}

\section{פעולות מנע ותוכניות מגירה}\label{sec:mitigation}

טבלה~\ref{tab:mitigation} מגדירה לכל סיכון פעולת מנע, תוכנית מגירה וציון שיורי אחרי פעולת המנע.

{\small
\begin{longtable}{|l|>{\raggedright\arraybackslash}p{6.4cm}|>{\raggedright\arraybackslash}p{4.4cm}|c|}
\caption{פעולות מנע, תוכניות מגירה וציון שיורי}\label{tab:mitigation}\\
\hline
\textbf{מזהה} & \textbf{פעולת מנע} & \textbf{תוכנית מגירה} & \textbf{שיורי} \\ \hline
\endfirsthead
\hline
\textbf{מזהה} & \textbf{פעולת מנע} & \textbf{תוכנית מגירה} & \textbf{שיורי} \\ \hline
\endhead
\en{DAT-01} & בחירת מקרים עם ציר זמן מתועד בכמה מקורות (\en{T1.1}); "מידע לא מספיק" הוא תוצאה נכונה שגם היא נבדקת & הצגת התוצאות כמקרי בוחן תיאוריים & 8 \\ \hline
\en{VAL-02} & הרצה לפי נקודת זמן, הנחיה שאוסרת ידע מחוץ למסמכים, ובדיקת תאריך כל מקור שצוטט & העדפת מקרים מאחרי תאריך הידע, ודיווח על המגבלה & 8 \\ \hline
\en{ALG-03} & טענה בלי מקור אינה נכנסת לציון; בדיקה אוטומטית שהקישור קיים והציטוט מופיע בו; השערים בקוד. הבדיקה אינה מבטיחה שהמסקנה נובעת מהמקור, ולכן הסיכון השיורי אינו מינימלי & ביקורת ידנית על מדגם טענות (\en{REQ-NFR-01}), והחלפת מודל & 10 \\ \hline
\en{SEC-04} & רק מידע ציבורי נכנס למערכת; שמות רכיבים כלליים בתיקוף; ארכיטקטורה שמאפשרת מודל מקומי (\en{REQ-CON-03}) & עצירת ההרצה, מחיקת המידע ודיווח למנחה & 5 \\ \hline
\en{LEG-05} & בדיקת תנאי השימוש לפני חיבור לכל מקור (\en{T1.2}); עדיפות לממשקי תכנות רשמיים & מקור חלופי מאותו סוג (טבלה~\ref{tab:data}) & 4 \\ \hline
\en{CST-06} & מדידת העלות לדוח באב-הטיפוס, ומודל קטן במשימות פשוטות, שבמאמר הנלווה כמעט לא נפל בביצועים \cite{zhao2025} & הגבלת מספר הספקים בהרצה & 2 \\ \hline
\en{ENT-07} & זיהוי הישות בקוד (\en{T2.1}) מול מרשם עם מזהה ייחודי, כולל חברות-בת ושמות קודמים, ואישור המשתמש לפני ההרצה & סימון "זיהוי לא ודאי" והעברה לבדיקה ידנית & 8 \\ \hline
\end{longtable}
}

אחרי פעולות המנע אף סיכון אינו נשאר קריטי. הגבוה ביותר הוא \en{ALG-03} (10), כי בדיקת הקישור אינה מבטיחה שהמסקנה נובעת מהמקור; אחריו \en{DAT-01}, \en{VAL-02} ו-\en{ENT-07} (8). \en{DAT-01} ו-\en{VAL-02} נוגעים לתוקף התיקוף, ולכן תוצאותיו יוצגו כמקרי בוחן ולא כהוכחה סטטיסטית.
